\chapter*{Summary}
This thesis studies the reliability of V2X communications in platooning scenarios, with the goal of understanding how to safely break up a platoon when link quality degrades. It builds on \emph{SafeSwitch}, a mechanism that links communication quality, controller selection, and inter-vehicle distance in order to avoid abrupt and potentially dangerous transitions.

The main contribution is the reconstruction of a multi-technology experimental environment built around OMNeT++, Veins, Plexe, SUMO, and Simu5G. Within this framework, the scenarios and metrics of the reference work were reproduced, the PDR monitoring logic and the \emph{SafeSwitch} state machine were reimplemented, and, most importantly, the cellular component was migrated from SimuLTE to Simu5G, updating the system to a more current 5G model.

The results show that the mechanism behaves consistently even on the new stack: in critical scenarios \emph{SafeSwitch} performs better than a \emph{naive} fallback, avoiding overly abrupt transitions and maintaining control of the platoon, while in the realistic scenario the migration to 5G introduces no functional regressions and, if anything, leads to fewer losses during cell handovers. The thesis therefore confirms that the architecture can be ported to a more modern toolchain and highlights how communication and control must be designed jointly to guarantee system safety.
